% =====================================================================
%  Enterprise Hybrid RAG — design and construction report
%
%  Build:  pdflatex report.tex && pdflatex report.tex
%  (twice, so the table of contents and cross-references resolve)
%
%  Every number in this document comes from an artefact committed in the
%  repository — data/evaluation/results.json, results_fallback.json,
%  scale_results.json — produced by a run that actually happened. None of
%  them is illustrative.
% =====================================================================
\documentclass[11pt,a4paper]{article}

\usepackage[T1]{fontenc}
\usepackage[utf8]{inputenc}
\usepackage{lmodern}
\usepackage[margin=2.4cm,bottom=2.8cm]{geometry}
\usepackage{microtype}
\usepackage{booktabs}
\usepackage{array}
\usepackage{xcolor}
\usepackage{enumitem}
\usepackage{titlesec}
\usepackage{fancyhdr}
\usepackage{tcolorbox}
\usepackage[hidelinks]{hyperref}

\definecolor{ink}{HTML}{1B1F24}
\definecolor{accent}{HTML}{0B5394}
\definecolor{rule}{HTML}{C9CFD6}
\definecolor{wash}{HTML}{F4F6F8}

\hypersetup{colorlinks=true, linkcolor=accent, urlcolor=accent, citecolor=accent}

\titleformat{\section}{\Large\bfseries\color{ink}}{\thesection}{0.6em}{}
\titleformat{\subsection}{\large\bfseries\color{ink}}{\thesubsection}{0.6em}{}
\titlespacing*{\section}{0pt}{2.0ex plus 1ex minus .2ex}{1.0ex plus .2ex}

\pagestyle{fancy}
\fancyhf{}
\renewcommand{\headrulewidth}{0.4pt}
\fancyhead[L]{\small\color{ink}Enterprise Hybrid RAG}
\fancyhead[R]{\small\color{ink}Design and construction report}
\fancyfoot[C]{\small\thepage}

\newcommand{\code}[1]{\texttt{\small #1}}
% T1 cannot stack a tilde on a circumflex; compose it explicitly so the
% name renders correctly under plain pdflatex.
\newcommand{\ecircumflextilde}{\~{\^e}}
\newcolumntype{R}{>{\raggedleft\arraybackslash}p{1.6cm}}

\newtcolorbox{finding}[1]{
  colback=wash, colframe=accent, boxrule=0.7pt, arc=2pt,
  left=8pt, right=8pt, top=6pt, bottom=6pt,
  title={\bfseries #1}, coltitle=white, fonttitle=\small}

\begin{document}

% ---------------------------------------------------------------- title
\begin{titlepage}
\centering
\vspace*{2.5cm}

{\Huge\bfseries\color{ink} Enterprise Hybrid RAG\par}
\vspace{0.8cm}
{\Large\color{ink} Design and construction report\par}
\vspace{0.4cm}
{\large\color{ink} Document-grounded question answering over PDF, DOCX and Markdown,\\
built without a retrieval framework\par}

\vspace{1.1cm}
{\color{ink}\normalsize T\'u Nguy\ecircumflextilde n\par}
\vspace{0.15cm}
{\color{ink}\small \today\par}

\vspace{1.4cm}
\begin{minipage}{0.86\textwidth}
\color{ink}\normalsize
\textbf{Abstract.} This report describes the design, construction and evaluation
of a retrieval-augmented question answering system built from primitives rather
than from a framework, so that every ranking decision is inspectable and every
published number is reproducible. The system fuses dense (bi-encoder) and sparse
(BM25) retrieval with Reciprocal Rank Fusion and reorders the result with a
cross-encoder, then answers only from the retrieved text and abstains when
retrieval confidence is too low.

The central empirical result is that the components are not independently good.
Replacing the offline fallback backends with real neural models makes dense
retrieval on its own \emph{worse} --- Recall@1 falls from 0.6550 to 0.5853 ---
while the full fused-and-reranked pipeline improves by 0.1512. The reason is
visible in a per-question-type breakdown: lexical matching is perfect on
reference codes and nearly useless on paraphrase; a sentence encoder is the
reverse. The value of hybrid retrieval is not an average of the two but the
recovery of both ceilings at once, and that recovery requires the cross-encoder
--- plain fusion scores \emph{worse} than either retriever on the category each
one owns.

A second experiment grows the corpus from 77 to 8548 chunks with in-domain
distractors while holding the questions fixed. Dense retrieval loses 35\% of its
Recall@1; the reranked pipeline loses 9\% and then stops, holding 0.7248 across a
twentyfold growth in the haystack. The reranker's advantage over dense retrieval
widens from $+0.2093$ to $+0.3450$: it does not merely survive scale, it is worth
more at scale.

The report also describes the measurement discipline that makes those numbers
worth reading --- guards that abort a run rather than publish a misleading
figure --- and three occasions on which that discipline caught the author's own
mistakes, including a continuous-integration check that reported success while
discarding the result it was meant to verify.
\end{minipage}

\vfill
{\color{ink}\small
Repository: \url{https://github.com/tu-h-nguyn/Enterprise-Hybrid-RAG}\\
All figures in this report are drawn from committed evaluation artefacts.\par}
\vspace{1cm}
\end{titlepage}

\tableofcontents
\newpage

% ================================================================ 1
\section{The problem, and the shape of the answer}

An organisation holds its operating knowledge in documents: a handbook, a
security policy, an incident runbook, a retention schedule. People do not want to
read them; they want to ask them a question. The failure mode of a language model
asked such a question directly is well known --- it produces a fluent answer with
no source, and the answer is sometimes wrong in a way the reader cannot detect.

The system described here constrains the answer to retrieved text, attaches a
citation to every claim, and refuses when retrieval is not confident enough. The
pipeline is deliberately ordinary:

\begin{center}
\small
\code{ingest} $\rightarrow$ \code{chunk} $\rightarrow$ \code{dense + BM25 index}
$\rightarrow$ \code{RRF fusion} $\rightarrow$ \code{cross-encoder rerank}
$\rightarrow$ \code{context build} $\rightarrow$ \code{grounded generation}
\end{center}

What is not ordinary is the decision to build it without a retrieval framework.
The core retrieval path has no LangChain, no LlamaIndex, no abstraction that
hides a ranking decision behind a configuration flag. That choice cost time and
bought two things that turned out to matter more than the time: every scoring
decision is a function that can be read in a minute and unit-tested against a
hand-computed example, and every number in the evaluation can be traced to the
code that produced it.

\subsection{Scope}

\begin{itemize}[leftmargin=1.2em,itemsep=2pt]
  \item \textbf{In scope.} Loading PDF, DOCX and Markdown with page and section
        provenance; chunking that preserves that provenance; two independent
        retrievers and an explicit fusion rule; reranking; an abstention gate;
        citation resolution; a FastAPI service; a Streamlit interface; an
        evaluation harness with a labelled dataset; continuous integration that
        regenerates the published numbers rather than trusting them.
  \item \textbf{Out of scope.} Multi-hop question decomposition, incremental
        index updates, authentication and multi-tenancy, and fine-tuning any
        model. Each is named in the limitations rather than quietly omitted.
\end{itemize}

% ================================================================ 2
\section{Architecture}

The system is a library with three entry points --- an HTTP API, a set of command
line scripts, and an evaluation harness --- all of which call the same method.
That is a deliberate constraint rather than an accident of layering: if the
evaluation called a reimplementation of retrieval, the published numbers would
describe code that no user ever runs.

\begin{center}
\small
\begin{tabular}{@{}lp{10.6cm}@{}}
\toprule
\textbf{Layer} & \textbf{Responsibility} \\
\midrule
\code{app/ingestion} & Loaders (PyMuPDF, python-docx, Markdown), cleaning, and a
metadata-aware chunker that carries page number and section heading onto every
chunk. Provenance is attached at load time and never re-derived, which is what
makes a citation trustworthy. \\
\addlinespace
\code{app/indexing} & Embedding backends, a Chroma or NumPy vector store, and a
persistent BM25 index. Each backend reports what it actually is, not what was
requested. \\
\addlinespace
\code{app/retrieval} & Dense retrieval, sparse retrieval, Reciprocal Rank Fusion,
rerankers, and a context builder that deduplicates and enforces a token budget. \\
\addlinespace
\code{app/generation} & Provider adapters (OpenAI-compatible, Anthropic, offline
extractive), prompt construction, and citation parsing that discards references
the model invented. \\
\addlinespace
\code{app/services} & \code{RagService} --- the single entry point --- plus the
abstention gate and document lifecycle. \\
\addlinespace
\code{app/evaluation} & Dataset schema, relevance resolver, retrieval and
generation metrics, and the experiment runners. \\
\bottomrule
\end{tabular}
\end{center}

\subsection{Backends degrade, and say so}

Every model-backed component can be set to \code{auto}. If the model hub is
unreachable --- which is the normal condition inside a restricted build
environment --- the embedder falls back to a TF-IDF plus truncated-SVD encoder,
the reranker falls back to a feature-based scorer, and generation falls back to
extractive sentence selection.

The important part is the second half of that sentence. Degrading silently would
be a liability: a benchmark run on fallbacks would publish numbers under a name
that implies neural models. So every artefact records which backend produced it,
\code{/health} reports the same, the report renderer emits an explicit warning
block, and the continuous integration job that publishes neural numbers
\emph{asserts} a 384-dimension sentence-transformer in the index manifest before
it measures anything.

% ================================================================ 3
\section{Design decisions worth defending}

\subsection{Reciprocal Rank Fusion, not score averaging}

Cosine similarity lives in $[-1, 1]$. BM25 is unbounded and corpus-dependent.
Averaging or weighting them makes the blend depend on one query's score
distribution, so a single outlier BM25 hit dominates the result. Reciprocal Rank
Fusion discards magnitudes and uses ranks only:
\[
  \mathrm{RRF}(d) \;=\; \sum_i \frac{w_i}{k + \mathrm{rank}_i(d)}, \qquad k = 60 .
\]
With $k = 60$, a document ranked first by one retriever cannot beat a document
ranked second or third by both --- agreement between two independent retrievers
is the signal being bought, and the constant is what makes it so. The
implementation is roughly fifteen lines and is unit-tested against a
hand-computed example, so the constant's effect is checkable rather than
folklore.

\subsection{Ground truth is answer spans, not chunk identifiers}

The evaluation dataset labels each question with its source documents and a
verbatim answer span, and resolves those to chunk identifiers at evaluation time.

Pinning chunk identifiers would have been easier and would have quietly destroyed
the most useful experiment in the project. Chunk identifiers depend on the
chunking configuration; the moment the chunk size changes, pinned labels address
nothing, and every metric collapses to zero in a way that looks exactly like a
retrieval failure. Because labels are spans, the chunk-size sweep in
Section~\ref{sec:chunk} is possible at all --- and, as measured in that section,
the resolved gold sets are identical in size and shape at both chunk sizes, which
is what makes the comparison fair.

\subsection{Confidence thresholds are per stage, because the scales differ}

Cosine similarity is in $[-1, 1]$; fused RRF scores are around $[0, 0.03]$; BM25
is unbounded; cross-encoder logits run roughly $[-11, +11]$; the lexical fallback
is around $[0, 5]$. A single ``confidence'' threshold across those scales is not a
tuning problem, it is a category error. Each stage therefore owns its own
threshold, and the abstention response records which one fired.

\subsection{The user interface is a client, not a second implementation}

The Streamlit interface holds no retrieval logic. It speaks HTTP to the same five
endpoints an external client would use, so what a visitor sees is what the API
serves.

That constraint was tested when the project needed to run on a single-process
host --- Streamlit Community Cloud and Hugging Face Spaces give you one process,
with nowhere to put a separate server. The obvious solution was to let the
interface import the service and call it directly. It was rejected for the reason
above. Instead the real ASGI application is started on a loopback socket inside
the same process and the interface keeps speaking HTTP to it. The transport
changes; the contract does not. A test asserts that every endpoint the interface
calls exists on the application that gets embedded, so the demonstration cannot
drift away from the deployment the documentation describes.

% ================================================================ 4
\section{Evaluation methodology}

\subsection{The dataset}

Fifty questions over twenty-two documents, written by hand against the corpus and
typed by the failure they are meant to expose:

\begin{center}
\small
\begin{tabular}{@{}lrp{8.6cm}@{}}
\toprule
\textbf{Type} & \textbf{n} & \textbf{What it tests} \\
\midrule
\code{factual}      & 14 & A fact stated once, in the words the question uses. \\
\code{keyword}      &  7 & A reference code --- \code{HB-7.2}, \code{FIN-TE-4}, \code{SEV1}. Lexical matching should win. \\
\code{paraphrased}  &  9 & The answer exists but shares no vocabulary with the question. Semantic matching should win. \\
\code{multi\_step}  &  7 & Evidence from two or more documents. \\
\code{terminology}  &  6 & A defined term asked plainly. \\
\code{unanswerable} &  7 & Nothing in the corpus answers it. The system should refuse. \\
\bottomrule
\end{tabular}
\end{center}

The seven unanswerable questions are not filler. Without them, an abstention gate
can be tuned to never abstain and the metrics improve.

\subsection{What makes a number worth publishing}

Three failures would each turn this evaluation into a confident lie, so each one
aborts the run rather than reporting a figure:

\begin{enumerate}[leftmargin=1.4em,itemsep=3pt]
  \item \textbf{A distractor that restates a labelled answer.} Relevance resolves
        per source document, so a distractor can never be \emph{labelled}
        relevant --- but it would be a correct answer scored as a miss. The
        corpus-scale run refuses to start if any distractor chunk contains a
        labelled span.
  \item \textbf{Gold chunks that move as the corpus grows.} If re-ingestion
        shifted the resolved gold set, a fall in Recall@1 would be a labelling
        artefact read as a scale effect. The run asserts the set is identical at
        every size --- per chunk size, because spans legitimately resolve to
        different chunks when the chunking changes.
  \item \textbf{An embedder whose width changes with the corpus.} The offline
        TF-IDF fallback derives its dimensionality from the corpus rank, so a
        sweep on it varies the encoder along with the variable under study.
        Continuous integration asserts a fixed 384 dimensions at every size
        before it will publish anything.
\end{enumerate}

A fourth discipline is procedural rather than assertive: quality metrics are
deterministic given the same corpus, dataset and backends, so a re-run is a
reproducibility test. The comparison script projects a result file down to its
deterministic parts, excludes wall-clock timings, and refuses outright to compare
runs made on different backends. Identical output means nothing drifted; a
difference is a real change worth committing.

\subsection{Where the numbers come from}

The development environment could not reach the model hub, so every neural figure
in this report was produced by a GitHub Actions workflow that pins the neural
backends explicitly rather than using \code{auto}. Under \code{auto}, an
unreachable hub would fall back silently and publish fallback numbers under a
workflow named ``neural'' --- precisely the failure the project is built to avoid.

% ================================================================ 5
\section{Results}

Unless stated otherwise: \code{all-MiniLM-L6-v2} as the bi-encoder,
\code{ms-marco-MiniLM-L-6-v2} as the cross-encoder, 22 documents and 77 chunks at
the shipped 256-token chunk size, 50 questions of which 43 are answerable.

\subsection{Retrieval comparison}

\begin{center}
\small
\begin{tabular}{@{}lrrrrrrr@{}}
\toprule
\textbf{Configuration} & \textbf{R@1} & \textbf{R@3} & \textbf{R@5} & \textbf{R@10} & \textbf{MRR} & \textbf{nDCG@5} & \textbf{ms} \\
\midrule
Dense only          & 0.5853 & 0.8256 & 0.8953 & 0.9302 & 0.7740 & 0.7869 & 16.58 \\
BM25 only           & 0.6434 & 0.8023 & 0.8023 & 0.8023 & 0.7597 & 0.7660 & \textbf{0.78} \\
Hybrid (RRF)        & 0.6434 & 0.8023 & 0.8488 & \textbf{1.0000} & 0.7910 & 0.7862 & 17.04 \\
\textbf{Hybrid + Reranker} & \textbf{0.7946} & \textbf{0.8953} & 0.9070 & \textbf{1.0000} & \textbf{0.9085} & \textbf{0.8937} & 699.95 \\
\bottomrule
\end{tabular}
\end{center}

Two columns carry more information than the headline. BM25's recall is
\emph{identical} at ranks 3, 5 and 10 --- it finds the chunk in the first three
results or it never finds it, which is what a lexical matcher does when the
query's words are not in the text. And both fused configurations reach 1.0000 at
rank 10: everything this corpus can answer is inside ten candidates, so from
there the problem is entirely ordering. That is why a reranker can reach 0.7946
at rank 1 without a better retriever underneath it.

\subsection{The components are not independently good}

\begin{center}
\small
\begin{tabular}{@{}lrrrr@{}}
\toprule
\textbf{Configuration} & \textbf{R@1 fallback} & \textbf{R@1 neural} & \textbf{MRR fallback} & \textbf{MRR neural} \\
\midrule
Dense only   & 0.6550 & 0.5853 & 0.7785 & 0.7740 \\
BM25 only    & 0.6434 & 0.6434 & 0.7597 & 0.7597 \\
Hybrid (RRF) & 0.6550 & 0.6434 & 0.7784 & 0.7910 \\
\textbf{Hybrid + Reranker} & 0.6434 & \textbf{0.7946} & 0.7541 & \textbf{0.9085} \\
\bottomrule
\end{tabular}
\end{center}

A better embedder made dense retrieval on its own \emph{worse}, left BM25 exactly
where it was --- as it must, since BM25 does not use the embedder, which is a
useful check that nothing else drifted between the two runs --- and moved plain
fusion by less than a question. Only the full pipeline gained, and it gained
$+0.1512$ Recall@1 and $+0.1544$ MRR.

The explanation is entirely in the breakdown by question type.

\begin{center}
\small
\begin{tabular}{@{}lrrrr@{}}
\toprule
\textbf{Recall@1} & \textbf{TF-IDF fallback} & \textbf{Dense (MiniLM)} & \textbf{BM25} & \textbf{Hybrid + Reranker} \\
\midrule
\code{keyword}     & \textbf{1.0000} & 0.4286 & \textbf{1.0000} & \textbf{1.0000} \\
\code{paraphrased} & 0.1111 & \textbf{0.4444} & 0.1111 & \textbf{0.4444} \\
\code{terminology} & 0.8333 & 0.3333 & 0.8333 & \textbf{1.0000} \\
\code{factual}     & 0.8571 & 0.9286 & 0.8571 & \textbf{1.0000} \\
\code{multi\_step} & 0.4524 & 0.4524 & 0.3810 & 0.4524 \\
\bottomrule
\end{tabular}
\end{center}

The TF-IDF fallback is perfect on keyword questions, because character n-grams
match a reference code literally, and nearly useless on paraphrase. MiniLM is the
mirror image. A sentence encoder compresses \code{FIN-TE-4} into a dense vector
and loses it.

\begin{finding}{The result the project exists to demonstrate}
Neither retriever is better. They fail in opposite directions. Fused and
reranked, the pipeline holds \emph{both} ceilings at once --- keyword 1.0000 and
paraphrased 0.4444, each matching the best single retriever --- and beats both of
them on terminology, 1.0000 against 0.8333. Its MRR is a perfect 1.0000 on four
of the five answerable categories.
\end{finding}

And fusion alone does not get there. Plain RRF scores 0.5714 on keyword and
0.2222 on paraphrased: \emph{worse than BM25 and worse than dense respectively},
on the category each of them owns. Averaging two rankings that disagree lands
between them. The cross-encoder, which reads query and chunk jointly rather than
comparing two independently compressed vectors, is what turns a fused candidate
list back into the best of both.

\subsection{Does it survive a bigger corpus?}
\label{sec:scale}

Seventy-seven chunks is a haybale, not a haystack: forty of them are the answer to
some question, and the top five cover 6.5\% of everything there is. Recall@5
saturates, and none of it answers whether the ranking survives a realistic index.

So the fifty questions, their answer spans, the chunking and every retrieval
setting were pinned, and only the corpus grew --- with distractor documents on the
\emph{same topics, in the same register, with the same policy vocabulary}, for
other fictional companies. Padding with off-topic prose would have proved
nothing: a policy question never retrieves a novel, recall would have stayed
flat, and the result would have been a rigged win.

\begin{center}
\small
\begin{tabular}{@{}rrrrrrr@{}}
\toprule
\textbf{Chunks} & \textbf{Docs} & \textbf{Gold share} & \textbf{Dense} & \textbf{BM25} & \textbf{Hybrid} & \textbf{Hybrid + Rerank} \\
\midrule
77   & 22   & 51.95\% & 0.5853 & 0.6434 & 0.6434 & \textbf{0.7946} \\
423  & 162  & 9.46\%  & 0.4690 & 0.6085 & 0.5853 & \textbf{0.7248} \\
1704 & 677  & 2.35\%  & 0.4341 & 0.5969 & 0.5620 & \textbf{0.7248} \\
8548 & 3422 & 0.47\%  & 0.3798 & 0.5969 & 0.4922 & \textbf{0.7248} \\
\midrule
\multicolumn{3}{@{}l}{\textbf{Change over 111$\times$}} & $-0.2055$ & $-0.0465$ & $-0.1512$ & $\mathbf{-0.0698}$ \\
\bottomrule
\end{tabular}
\end{center}

\begin{finding}{The reranker is worth more at scale, not less}
Dense retrieval gives up 35\% of its Recall@1 across the range. The reranked
pipeline gives up 9\% and then \emph{stops}: 0.7248 at 423 chunks, 0.7248 at
1704, 0.7248 at 8548 --- flat across a twentyfold growth, while the gold chunk
goes from one in 1.9 to one in 213. Its advantage over dense alone widens from
$+0.2093$ to $+0.3450$. That is the argument for paying its latency, and it is
completely invisible at the size the rest of the evaluation runs at.
\end{finding}

BM25 is nearly flat ($-0.0465$, two questions), which is the expected shape:
exact lexical matching does not care how much other text exists, only whether
something else matches better.

\paragraph{Are the distractors actually hard?} Asserting it would be worthless,
so every row measures it. At 8548 chunks, a distractor takes rank 1 on 41.9\% of
answerable questions under dense retrieval and distractors hold 75.4\% of the top
five. They compete. The reranked pipeline holds its own rate at a flat 0.1628
from 423 chunks upward --- the same shape as its Recall@1.

\subsection{Chunk size, and changing a default on evidence}
\label{sec:chunk}

The chunk-size ablation had said for some time that 256 tokens beat the shipped
500 by 9.3 points of Recall@1 for the production configuration. The default was
left alone, and the report said why: one fifty-question corpus of 44 chunks is
thin evidence, and a chunk size that wins on a corpus that small may be a fact
about the corpus rather than about chunking.

That objection was answerable, so it was answered. The scale sweep runs the whole
ladder at both sizes, matched by how many distractor \emph{documents} were added,
because the same documents become a different number of chunks at each size.

\begin{center}
\small
\begin{tabular}{@{}rrrrrrr@{}}
\toprule
& \multicolumn{3}{c}{\textbf{256 tokens}} & \multicolumn{3}{c}{\textbf{500 tokens}} \\
\cmidrule(lr){2-4}\cmidrule(lr){5-7}
\textbf{Distractor docs} & \textbf{chunks} & \textbf{R@1} & \textbf{MRR} & \textbf{chunks} & \textbf{R@1} & \textbf{MRR} \\
\midrule
0    & 77   & \textbf{0.7946} & \textbf{0.9085} & 44   & 0.7016 & 0.8568 \\
140  & 423  & \textbf{0.7248} & \textbf{0.8384} & 235  & 0.6550 & 0.8101 \\
655  & 1704 & \textbf{0.7248} & \textbf{0.8289} & 936  & 0.6318 & 0.7950 \\
3400 & 8548 & \textbf{0.7248} & \textbf{0.8285} & 4658 & 0.6318 & 0.7868 \\
\bottomrule
\end{tabular}
\end{center}

256 wins Recall@1 at every size, by 0.0930 at the largest, and the gap does not
shrink as the corpus grows. It also halves the cost of the dominant stage:
685.39\,ms against 1370.06\,ms in the same sweep, because the cross-encoder scores
passages half as long. These rows then reproduced exactly on a second runner,
with the roles of default and alternative swapped between the two runs. The
default moved to 256 on that evidence.

\paragraph{The cost, stated rather than buried.} 500 wins Recall@5, 0.8566
against 0.8062 at the largest corpus. Two things decide it anyway. First, the
resolved gold sets are identical at both chunk sizes --- mean 1.209 relevant
chunks per question, distribution $\{1{:}36,\,2{:}5,\,3{:}2\}$ --- so this is not
a labelling artefact; that was checked, not assumed. Second, Recall@5 is not a
like-for-like comparison across chunk sizes: five chunks of 500 tokens hand the
generator twice the text that five chunks of 256 do. The metrics where a longer
chunk is \emph{structurally} advantaged are Recall@1 and MRR --- a bigger chunk
is a bigger target to hit --- and 256 wins those anyway, by nine points. The
metric 500 wins is the one where it is given twice as much room. A comparison at
a fixed context budget, $\mathrm{top\_k}=10$ at 256 against $\mathrm{top\_k}=5$ at
500, is recorded as the follow-up this does not attempt.

\subsection{What approximate search costs}

Every size was also run against both vector stores --- Chroma's HNSW index and
exhaustive cosine --- because that difference is normally assumed rather than
measured. With \code{all-MiniLM-L6-v2}, Recall@1 is identical for every
configuration at every size, and the single difference anywhere in the sweep is
0.0023 of dense MRR at the largest corpus. BM25 never touches the vector store
and reads identical everywhere, which is the control confirming nothing else
differed between the two runs.

That result does not transfer to any embedding. Under the offline TF-IDF
fallback the same comparison loses 0.0233 \emph{Recall@1} to HNSW and is not
reproducible at all: repeated runs of the same command returned dense Recall@1 of
0.3643, 0.3411, 0.3876 and 0.3876 --- the last two with BLAS threading pinned to
one core, which ruled out floating-point reduction order and pointed at the index
itself. Exhaustive search returned 0.4109 twice and matched on all 86 compared
values. ``HNSW is fine'' is a statement about how well separated these embeddings
are, not about HNSW.

\subsection{Cost}

\begin{center}
\small
\begin{tabular}{@{}lrr@{}}
\toprule
& \textbf{77 chunks} & \textbf{8548 chunks} \\
\midrule
Index build                & $\sim$9\,s        & $\sim$3.7\,min \\
Dense query (exhaustive)   & $\sim$15\,ms      & $\sim$21\,ms \\
BM25 query                 & under 1\,ms       & $\sim$16\,ms \\
Hybrid + rerank query      & $\sim$0.7\,s      & $\sim$0.7\,s \\
\bottomrule
\end{tabular}
\end{center}

Rounded on purpose: these are wall-clock figures on a shared runner and move by
tens of percent between runners. The shape is the claim. Reranking is flat
because it always rescores a fixed top-$k$ --- the cross-encoder never sees the
corpus --- and it is 98\% of query latency at 77 chunks and still 94\% at 8548.
Growing the corpus by $111\times$ does not change what dominates a request; what
it changes is the index build, which is where the corpus is actually paid for.

\subsection{Generation}

Generation is measured on the \emph{extractive} backend, because continuous
integration has no API key. These numbers describe sentence selection, not a
language model, and are reported for completeness rather than as a claim about
answer quality.

\begin{center}
\small
\begin{tabular}{@{}lr@{}}
\toprule
\textbf{Metric} & \textbf{Value} \\
\midrule
groundedness (answer bigrams present in context) & 0.6434 \\
answer\_f1 (SQuAD-style token F1)                 & 0.3137 \\
span\_coverage (labelled span appears in answer)  & 0.5078 \\
citation precision                                & 0.6512 \\
\textbf{correct refusal} (unanswerable)           & \textbf{0.8571} \\
\textbf{over-refusal} (answerable)                & \textbf{0.3488} \\
\bottomrule
\end{tabular}
\end{center}

The gate correctly refuses six of seven unanswerable questions and also refuses
fifteen of forty-three answerable ones. Both rates are unchanged from the
fallback run and did not move when the chunk size changed, which is itself
informative: the gate's behaviour here is dominated by the extractive backend's
exact-token matching rather than by retrieval quality. Measuring it properly
requires a real language model, and that is the first open item.

% ================================================================ 6
\section{How the system is used}

\subsection{Three surfaces, one implementation}

\paragraph{As an HTTP service.} The intended production shape. \code{POST /query}
takes a question and returns an answer, citations resolved to document, page and
section, the retrieved chunks, and a trace of what each stage did and how long it
took. \code{GET /health} reports index readiness and --- importantly --- which
backends were actually selected, so a degraded deployment is visible from
outside. \code{POST /documents/upload} and \code{POST /documents/index} cover the
ingestion lifecycle.

\paragraph{As a user interface.} A Streamlit application for the people who will
not write \code{curl}. It shows the answer, every citation as an expandable block
with its snippet, the retrieved chunks with their component scores, and a debug
panel containing the raw trace. When the abstention gate fires, it shows which
threshold fired rather than a generic apology.

With \code{API\_URL} pointing at a running service, it is a thin client. With
\code{API\_URL} unset, it starts the service inside its own process and builds the
index on first visit --- which is what makes it deployable to a single-process
host with no configuration at all.

\paragraph{As a library and a set of scripts.} \code{scripts/ingest.py} builds
the indexes; \code{scripts/evaluate.py} scores retrieval or generation;
\code{scripts/benchmark.py} runs the full comparison and ablation and writes
\code{results.json} and \code{report.md}; \code{scripts/scale\_experiment.py} runs
the corpus-scale sweep; \code{scripts/compare\_results.py} decides whether two
runs agree on their deterministic parts.

\subsection{Getting it running}

\begin{center}
\small
\begin{tabular}{@{}p{5.4cm}p{9.6cm}@{}}
\toprule
\textbf{Goal} & \textbf{Command} \\
\midrule
Local service and interface &
\code{pip install -r requirements.txt}\newline
\code{python scripts/ingest.py}\newline
\code{uvicorn app.main:app --reload}\newline
\code{API\_URL=http://localhost:8000 streamlit run frontend/streamlit\_app.py} \\
\addlinespace
Everything in one process &
\code{streamlit run frontend/streamlit\_app.py} \newline (no \code{API\_URL}: the
service starts in-process and the index is built on first visit) \\
\addlinespace
Containers &
\code{docker compose up --build} \newline (API on 8000, interface on 8501) \\
\addlinespace
Reproduce the numbers &
\code{python scripts/benchmark.py} \newline
\code{python scripts/scale\_experiment.py --extra-chunk-sizes 500} \\
\bottomrule
\end{tabular}
\end{center}

\subsection{Choosing a generation backend}

The embedder and the reranker are public models and need no token of any kind.
Only generation needs a key, and only if a generative answer is wanted:

\begin{itemize}[leftmargin=1.2em,itemsep=2pt]
  \item \textbf{Nothing configured} --- the extractive backend selects supporting
        sentences from the retrieved context. No key, no network, and every
        artefact says so.
  \item \textbf{A local model server} --- point \code{LLM\_BASE\_URL} at an
        OpenAI-compatible endpoint such as Ollama. No key leaves the machine and
        no document does either. Two traps are documented because they cost time:
        \code{LLM\_API\_KEY} must be non-empty even though such a server ignores
        it, and \code{localhost} inside a container must become
        \code{host.docker.internal}.
  \item \textbf{A hosted provider} --- OpenAI-compatible or Anthropic, selected by
        environment variable. Both adapters are covered by tests that stand up a
        real HTTP server on a loopback port and assert what each one actually
        sends: the endpoint, the authentication header, and the request shape.
\end{itemize}

No secret is ever written to a file in the repository, baked into an image, or
printed in a log.

\subsection{Where it fits, and where it does not}

The system suits an internal corpus of policy-like documents --- handbooks,
runbooks, contracts, standards --- where answers must be attributable and a
confident wrong answer is worse than an admission of ignorance. It does not suit
open-domain question answering, questions requiring synthesis across many
documents (see \code{multi\_step} in the results), or a corpus that changes
continuously, because indexing is a full rebuild rather than an incremental
update.

% ================================================================ 7
\section{Engineering process}

\subsection{What continuous integration actually enforces}

\begin{itemize}[leftmargin=1.2em,itemsep=3pt]
  \item \textbf{Lint, types and tests} --- \code{ruff}, \code{mypy} over
        \code{app/}, \code{scripts/} and \code{frontend/}, and 151 tests behind a
        coverage floor that fails the build if it drops.
  \item \textbf{A blocking dataset audit.} Evaluation labels are resolved against
        freshly ingested chunks; if any label stops resolving, the build fails.
        This gate was verified by deliberately corrupting an answer span and
        confirming a non-zero exit --- a gate nobody has seen fail is a gate
        nobody should trust.
  \item \textbf{A container that is built, ingested into, served and queried.}
        Not built and declared working.
  \item \textbf{An image-size ceiling and a \code{+cpu} assertion}, because torch
        arrives as a transitive dependency and the Linux wheel on the public
        index is the CUDA build --- roughly 3\,GB of wheels for a container with
        no GPU.
  \item \textbf{Regenerated benchmarks.} The neural numbers are recomputed and
        compared against the committed ones on quality metrics only, since
        latency is wall-clock. Identical output means the benchmark reproduced;
        it commits a new artefact only when a quality metric actually moved.
\end{itemize}

\subsection{Three times the discipline caught the author}

A report that lists only successes is not evidence of engineering judgement. The
following are the three most useful failures of this project, all of them found
by machinery rather than by inspection.

\begin{finding}{1. A benchmark that had never run}
The retrieval evaluator unpacked \code{retrieve()} as
\code{\_, final, \_} when the method returns
\code{(candidates, trace, stopwatch)} --- it was iterating a trace object. The
phase was marked complete in the handover document. It had never been executed
end to end. The fix was one line; the lesson was that ``done'' in a status table
is a claim, not a measurement.
\end{finding}

\begin{finding}{2. Numbers transcribed into prose go stale}
Wall-clock figures were copied into the documentation --- \code{1416.77\,ms},
\code{154.99\,s} --- and a later run on a different runner invalidated all of
them. The mistake was the transcription, not the figures: shared-runner timings
move by tens of percent and say nothing the shape does not. They are now rounded,
labelled as one run's wall-clock, and pointed at the committed artefact. A script
checks that every four-decimal number in both READMEs traces back to a committed
artefact; the only figures that do not are a Python version, an image size
explicitly labelled as the \emph{before}, and an example payload labelled with
the default it was captured at.
\end{finding}

\begin{finding}{3. A check that reported success while it was wrong}
Chroma's HNSW index is approximate, and two identical runs disagreed on a handful
of deep-rank values. To stop that rewriting the artefact on every run, the
comparison script gained an \code{--allow-drift} prefix. Then the chunk size
changed, which swapped which sizes each store runs at --- so the regenerated
artefact carried \code{chroma/c256/...} rows where the committed one had
\code{chroma/c500/...}. Every one of those keys existed on exactly one side,
every one began with \code{scale.chroma/}, and the guard counted them all as
tolerable noise. The job printed \emph{``reproduced the committed numbers
exactly; nothing to commit''} and discarded a correct measurement.

Drift now means a value moved. A key present on one side and absent on the other
is a structural change and no prefix excuses it. The test that reproduces the
original failure fails without the fix.
\end{finding}

The third is the one worth dwelling on. The guard was added for a real reason and
made the system quieter; it also made it lie, and it lied in the direction of
reporting success. A check that can only fail loudly is worth more than a check
that is usually right.

% ================================================================ 8
\section{Limitations}

Each of these bounds how far the numbers above generalise.

\begin{enumerate}[leftmargin=1.4em,itemsep=3pt]
  \item \textbf{The corpus is synthetic, and the labelled part of it is small.}
        Twenty-two documents, 77 chunks, 43 scored questions --- where one
        question is 0.0233. The per-type rows rest on six to fourteen questions
        each. The scale experiment grows the index to 8548 chunks, but with
        generated documents: hard negatives by construction, where a real corpus
        that size would hold both easier negatives and harder ones, such as
        near-duplicate revisions of the same policy. Treat every comparison as
        directional.
  \item \textbf{Generation is not measured with a real language model.} The
        retrieval numbers are real; the generation numbers describe sentence
        selection. Everything about the abstention gate is therefore provisional.
  \item \textbf{The abstention threshold is tuned on one corpus} and over-refuses
        34.88\% of answerable questions. That is the system's worst
        product-level flaw, and it is downstream of the previous item.
  \item \textbf{Token counts are a \code{chars/4} heuristic}, not a real
        tokenizer, so chunk sizes and context budgets are approximate and a real
        tokenizer would shift chunk boundaries.
  \item \textbf{Reranking dominates latency} --- 699.95\,ms against BM25's
        0.78\,ms. Any deployment must decide whether that precision is worth
        nearly three orders of magnitude, or whether to rerank only when the
        fusion margin is narrow.
  \item \textbf{Indexing is a full rebuild}, not an incremental upsert. Correct
        at this size, wrong at scale.
  \item \textbf{\code{multi\_step} is the weakest answerable category} (0.4524
        Recall@1, against a perfect 1.0000 MRR --- the \emph{first} relevant
        chunk is always ranked first and what is missing is the second one).
        Nothing here decomposes a question or retrieves iteratively, and no
        amount of reranking substitutes for that.
\end{enumerate}

% ================================================================ 9
\section{What comes next, in order}

\begin{enumerate}[leftmargin=1.4em,itemsep=3pt]
  \item \textbf{Measure generation against a real model.} A local server needs no
        key and keeps the corpus on the machine, and it unlocks the judge that
        refuses to run on the extractive backend. Items 2 and 3 are blocked on
        this, directly or indirectly.
  \item \textbf{Compare chunk sizes at a fixed context budget}
        ($\mathrm{top\_k}=10$ at 256 against $\mathrm{top\_k}=5$ at 500), which is
        the one comparison the chunk-size decision did not control for.
  \item \textbf{Re-tune the abstention thresholds} once there is a real generator
        to tune them against.
  \item \textbf{Attack \code{multi\_step}} with query decomposition or iterative
        retrieval --- the only weak category that reranking provably cannot fix.
  \item \textbf{Incremental indexing}, so a changed document does not require a
        full rebuild.
\end{enumerate}

% ================================================================ A
\appendix
\section{Repository map}

\begin{center}
\small
\begin{tabular}{@{}lp{10.2cm}@{}}
\toprule
\textbf{Path} & \\
\midrule
\code{app/ingestion/}  & Loaders, cleaner, metadata-aware chunker \\
\code{app/indexing/}   & Embedders, Chroma and NumPy vector stores, BM25 \\
\code{app/retrieval/}  & Dense, sparse, RRF fusion, rerankers, context builder \\
\code{app/generation/} & Providers, prompts, citation resolution \\
\code{app/evaluation/} & Dataset schema, resolver, metrics, experiment runners \\
\code{app/services/}   & RAG pipeline, abstention gate, document lifecycle \\
\code{frontend/}       & Streamlit interface and the in-process API starter \\
\code{scripts/}        & Ingest, evaluate, benchmark, corpus-scale experiment \\
\code{tests/}          & 151 unit and integration tests \\
\code{data/evaluation/}& The labelled dataset and every committed result artefact \\
\code{HANDOFF.md}      & Interface contracts, status by phase, and what is open \\
\bottomrule
\end{tabular}
\end{center}

\section{Reproducing every figure in this report}

\begin{enumerate}[leftmargin=1.4em,itemsep=2pt]
  \item \code{pip install -r requirements.txt}
  \item \code{python scripts/make\_demo\_corpus.py} --- writes the 22 documents
  \item \code{python scripts/ingest.py} --- 22 documents, 77 chunks
  \item \code{python scripts/benchmark.py} --- the comparison, the ablation and
        generation; writes \code{data/evaluation/results.json}
  \item \code{python scripts/scale\_experiment.py --extra-chunk-sizes 500} ---
        the corpus-scale sweep; writes \code{scale\_results.json}
\end{enumerate}

Without access to the model hub, steps 4 and 5 will run on the offline fallbacks
and every artefact they write will say so in its \code{backends} block and in the
warning at the top of its report. The neural figures in this document come from
the \code{benchmark.yml} workflow, which pins the backends and asserts a
384-dimension embedder before measuring.

\vspace{1.2em}
\noindent\rule{\textwidth}{0.4pt}
\vspace{0.4em}

\noindent\small\textit{Every figure in this report was taken from a committed
evaluation artefact produced by a run that actually executed. Where a number is
historical --- a previous image size, an example payload captured at an earlier
default --- it is labelled as such in the text.}

\end{document}
